\section*{3 Pseudorandomness Preliminaries}
Having completed our analysis of subset-wise symmetric polynomials, we now move on to setting the groundwork for our PRG construction and analysis.
\subsection*{3.1 Probability Basics}
Let $U_n$ denote the uniform distribution over $\{0,1\}^n$. We will simply write $U$ if $n$ is clear from context. For $f:\{0,1\}^n\to\mathbb R$, as a shorthand, we write $\mathbb E[f]$ to denote $\mathbb E[f(U)]$ and $\operatorname{Var}[f]$ to denote $\operatorname{Var}[f(U)]$. If $X$ is a distribution over $\{0,1\}^n$, we say that $X$ $\epsilon$-fools $f$, or $X$ fools $f$ with error $\epsilon$, if $|\mathbb E[f(X)]-\mathbb E[f]|\le\epsilon$. We say that $X$ $\epsilon$-fools a family $\mathcal F$ of functions, if it $\epsilon$-fools every $f\in\mathcal F$.
\subsection*{3.2 Small Bias}
A parity function is a function of the form $f(x)=\bigoplus_{i\in I}x_i$ for some set $I\subseteq[n]$. We say that a random variable $Y\in\{0,1\}^n$ is $\delta$-biased if it $\delta$-fools all parity functions. We say that $Y$ is $n'$-wise $\delta$-biased if it $\delta$-fools all parity functions on at most $n'$ bits, i.e., all parity functions with $|I|\le n'$. There are explicit constructions of $n'$-wise $\delta$-biased distributions that can be sampled with $O(\log(n'/\delta)+\log\log n)$ truly random bits [25,3].
Recall that for a function $f:\{0,1\}^n\to\mathbb R$ with Fourier expansion $f=\sum_{S\subseteq[n]}\widehat f(S)\cdot\chi_S$, the $L_1$ norm of $f$ is defined by $L_1(f)=\sum_{S\subseteq[n]}|\widehat f(S)|$. This norm is subadditive ($L_1(f+g)\le L_1(f)+L_1(g)$) and submultiplicative ($L_1(f\cdot g)\le L_1(f)\cdot L_1(g)$). Functions with bounded $L_1$ norm are fooled by small-bias distributions:
\begin{authorclaim}{Claim5}
If $f:\{0,1\}^n\to\mathbb R$ and $Y$ is $\delta$-biased, then $Y$ fools $f$ with error $2\delta\cdot L_1(f)$.
\end{authorclaim}
We will also rely on the following "XOR lemma" for small-bias distributions.
\begin{authorclaim}{Lemma6 ([16,24])}
Let $0<\delta<\epsilon\le1$. Let $f_1,\ldots,f_k:\{0,1\}^n\to[-1,1]$ depend on disjoint variable sets, and define $f(x)=\prod_{i=1}^kf_i(x)$. If every $\delta$-biased distribution $\epsilon$-fools every $f_i$, then every $\delta^k$-biased distribution fools $f$ with error $16^k\cdot2\epsilon$.
\end{authorclaim}
\subsection*{3.3 Limited Independence}
For $p\in[0,1]$, let $\operatorname{Bernoulli}(p)^{\otimes n}$ denote the distribution over $\{0,1\}^n$ where the bits are i.i.d. and each bit has expectation $p$. For example, $U_n=\operatorname{Bernoulli}(1/2)^{\otimes n}$. For a set $I=\{i_1<i_2<\cdots<i_\ell\}\subseteq[n]$ and a string $z\in\{0,1\}^n$, we let $z|_I=z_{i_1}z_{i_2}\cdots z_{i_\ell}\in\{0,1\}^\ell$. We say that $Z\in\{0,1\}^n$ is $\gamma$-almost $k$-wise independent with marginals $p$ if for every set $I\subseteq[n]$ with $|I|\le k$, the total variation distance between $Z|_I$ and $\operatorname{Bernoulli}(p)^{\otimes|I|}$ is at most $\gamma$.
\begin{authorclaim}{Claim7}
For every $n,k,C\in\mathbb N$ and $\gamma>0$, there is an explicit $\gamma$-almost $k$-wise independent distribution with marginals $p=1-2^{-C}$ that can be sampled with $O(Ck+\log(1/\gamma)+\log\log n)$ truly random bits.
\end{authorclaim}
\begin{proof}
Sample $Y\in\{0,1\}^{Cn}$ from a $(Ck)$-wise $(2^{-Ck/2-1}\gamma)$-biased distribution. Note that as discussed above $Y$ can be sampled using $O(\log(2^{Ck}/\gamma)+\log\log n)=O(Ck+\log(1/\gamma)+\log\log n)$ truly random bits. Divide $Y$ into $n$ blocks $Y^{(1)},\ldots,Y^{(n)}\in\{0,1\}^C$, and set $Z_i=0\Longleftrightarrow Y^{(i)}=1^C$. The desired distribution is $Z\in\{0,1\}^n$.
To prove correctness, let $f:\{0,1\}^n\to\{0,1\}$ be any test function depending on only $k$ variables. There is a function $g:\{0,1\}^{Cn}\to\{0,1\}$ depending on only $Ck$ variables such that $f(Z)=g(Y)$. By Claim5,
\[
|\mathbb E[f(Z)]-\mathbb E[f(\operatorname{Bernoulli}(p)^{\otimes n})]|
=|\mathbb E[g(Y)]-\mathbb E[g]|
\le2^{-Ck/2-1}\cdot2\gamma\cdot L_1(g)\le\gamma.
\]
\end{proof}
The expectation parameter $p$ can be "amplified" by drawing independent samples and combining with a coordinate-wise conjunction:
\begin{authorclaim}{Claim8}
Let $Z$ be $\gamma$-almost $k$-wise independent with marginals $p$. Draw $t$ independent samples $z^{(1)},\ldots,z^{(t)}\sim Z$, and let $Z'=z^{(1)}\wedge\cdots\wedge z^{(t)}$. Then $Z'$ is $(t\gamma)$-almost $k$-wise independent with marginals $p^t$.
\end{authorclaim}
\begin{proof}[Proof sketch]
The proof is a simple hybrid argument. Draw $t$ independent samples $r^{(1)},\ldots,r^{(t)}\sim\operatorname{Bernoulli}(p)^{\otimes n}$, and let $Z^{(i)}=z^{(1)}\wedge\cdots\wedge z^{(i)}\wedge r^{(i+1)}\wedge\cdots\wedge r^{(t)}$. One can show by induction on $i$ that $Z^{(i)}$ is $(i\gamma)$-almost $k$-wise independent with marginals $p^t$.
\end{proof}
\subsection*{3.4 PARITY $\circ$ AND Formulas}
Recall that our main result (Theorem1) is a PRG for read-once depth-2 $\mathrm{AC}^0[\oplus]$. For most of the paper, we will focus on the special case that the root gate is $\oplus$ and its immediate children are $\wedge$ gates. That is, define a PARITY $\circ$ AND formula to be a function of the form $f(x)=\bigoplus_{i=1}^mf_i(x)$, where each $f_i$ is a conjunction of literals, i.e., variables or their negations. We refer to $f_1,\ldots,f_m$ as the terms of $f$. We say that the formula is read-once if each variable $x_i$ appears in at most one term. Most of our effort will be spent fooling read-once PARITY $\circ$ AND formulas. Note that this is a slight generalization of read-once $\mathbb F_2$-polynomials due to the availability of $\neg$ gates. We will explain in Section5.6 why it is sufficient to focus on this special case.
The width of a term is the number of variables in the term; the width of $f$ is the maximum width of its terms. The length of $f$ is $m$, the number of its terms.
For convenience, if $f$ is a function taking values in $\{0,1\}$, we let $\mathfrak f=(-1)^f$. That way, if $f$ is a PARITY $\circ$ AND formula, $\mathfrak f=\prod_{i=1}^m\mathfrak f_i$.
\subsection*{3.5 Restrictions}
A restriction is a string $x\in\{0,1,?\}^n$; intuitively, $x_i=?$ means that $x_i$ has still not been assigned a value. We define an associative composition operation on restrictions by the formula $(x\circ x')_i=x_i$ if $x_i\ne?$, and $x'_i$ otherwise.
For a function $f$ on $\{0,1\}^n$, the restricted function $f|_x$ on $\{0,1\}^n$ is defined by $f|_x(x')=f(x\circ x')$.
A restriction $x$ can be specified by two strings $y,z\in\{0,1\}^n$ using the following notation.\footnote{With apologies, we here flip the order of the arguments to Res compared to the notation used in the authors' prior work [13].} Define $\operatorname{Res}:\{0,1\}^n\times\{0,1\}^n\to\{0,1,?\}^n$ by $(\operatorname{Res}(y,z))_i=?$ if $z_i=1$, and $y_i$ if $z_i=0$. In words, $z$ indicates the ? positions, and $y$ provides the bits in the non-? positions.
\subsection*{3.6 Pseudorandom Restrictions}
Let $Y,Z$ be distributions over $\{0,1\}^n$, and let $X=\operatorname{Res}(Y,Z)$. For a function $f:\{0,1\}^n\to\mathbb R$, we say that the distribution $X$ preserves the expectation of $f$ with error $\epsilon$ if $|\mathbb E[f|_X(U)]-\mathbb E[f]|\le\epsilon$.
An equivalent condition is that $|\mathbb E[f(Y+Z\wedge U)]-\mathbb E[f]|\le\epsilon$, where $+$ denotes addition over $\mathbb F_2^n$ and $\wedge$ denotes coordinate-wise conjunction. This second condition is the "pseudorandomness plus noise" perspective [18] (the string $Z\wedge U$ can be thought of as a noise vector.)
If $f$ takes on values in $\{0,1\}$, for each particular value $z$ that $Z$ might take on, we define the bias function [16] $\widetilde f_z:\{0,1\}^n\to[-1,1]$ by $\widetilde f_z(x)=\mathbb E[\mathfrak f(x+z\wedge U)]$. (We use $\mathfrak f$ rather than $f$ simply for convenience.) The statement that $X$ preserves the expectation of $f$ with error $\epsilon$ is also equivalent to the condition
\[
\left|\mathbb E_Z[\mathbb E_Y[\widetilde f_Z(Y)]-\mathbb E[\mathfrak f]]\right|\le2\epsilon.
\]
When $z$ is clear from context, we will just write $\widetilde f$ instead of $\widetilde f_z$.
If $X$ is a distribution over $\{0,1,?\}^n$ and $t\in\mathbb N$, let $X^{\circ t}$ denote the distribution over $x\in\{0,1,?\}^n$ obtained by drawing independent samples $x^{(1)},\ldots,x^{(t)}\sim X$ and composing them, $x=x^{(1)}\circ\cdots\circ x^{(t)}$.
Suppose $\mathcal F$ is a class of Boolean functions that is closed under restriction. If $X$ preserves the expectation of every $f\in\mathcal F$ with error $\epsilon$, then $X^{\circ t}$ preserves the expectation of every $f\in\mathcal F$ with error $t\epsilon$. Furthermore, informally, if $X$ "has ?-probability $p$", then $X^{\circ t}$ "has ?-probability $p^t$". To be precise, we can consider the case $X=\operatorname{Res}(Y,Z)$ where $Z$ is $\gamma$-almost $k$-wise independent with marginals $p$. Then the distribution of ? positions in $X^{\circ t}$ is described by Claim8.
